\rok{Новембар 2023.}{E}

Први практични тест из Алгоритама и структура података

\zadatak{1}{Брисање елемената из стека (Stack Pop)}
Написати методу која имплементира стандардан начин за брисање елемената из стека.

\textbf{Додатне напомене за решавање задатка:}
\begin{itemize}
    \item Алгоритам написати у оквиру закоментарисане методе:
    \begin{Verbatim}[fontsize=\footnotesize, numbers=left, xleftmargin=10mm]
pop?(int value)
    \end{Verbatim}
    \item Поменута метода налази се у оквиру класе \texttt{Stack} која се налази у придруженом шаблону.
    \item Обезбедити код у \texttt{Main} методи за тестирање алгоритма.
\end{itemize}



\zadatak{2}{Елиминисање негативних бројева из низа}
Креирати функцију која ће из статичког низа елиминисати све негативне целе бројеве. Алгоритам је неопходно написати са линеарном временском комплексношћу $O(n)$.

\textbf{Додатни захтеви за решавање задатка:}
\begin{itemize}
    \item Алгоритам мора резултовати креирањем новог статичког низа.
    \item Захтеван потпис методе:
    \begin{Verbatim}[fontsize=\footnotesize, numbers=left, xleftmargin=10mm]
public static long[] EliminateNegatives (long[] arr)
    \end{Verbatim}
\end{itemize}

\textbf{Примери:}
\begin{itemize}
    \item \textbf{Улаз:} \texttt{[2,-3,5,-2,10]} \\
          \textbf{Излаз:} \texttt{[2,5,10]}
    \item \textbf{Улаз:} \texttt{[]} \\
          \textbf{Излаз:} \texttt{[]}
\end{itemize}



\zadatak{3}{Кришнамуртијев број}
Креирати алгоритам који проверава да ли се од елемената циркуларне листе може креирати Кришнамуртијев број. Кришнамуртијев број је број чија сума факторијела сваке цифре је једнака самом броју. Алгоритам је неопходно написати са линеарном временском комплексношћу $O(n)$.

\textbf{Додатни захтеви за решавање задатка:}
\begin{itemize}
    \item Неопходно је имплементирати функцију \texttt{Factorial} која треба да пронађе факторијел прослеђеног броја.
    \item Број 145 је Кришнамуртијев број: $1! + 4! + 5! = 1 + 24 + 120 = 145$
    \item Број 15 није Кришнамуртијев број: $1! + 5! = 1 + 120 = 121$
    \item Задатак решавати помоћу структуре циркуларне листе (класа \texttt{CircularList}) која је дата у оквиру шаблона задатка.
    \item Захтеван потпис методе:
    \begin{Verbatim}[fontsize=\footnotesize, numbers=left, xleftmargin=10mm]
public bool IsListKrishnamurthyNumber()
    \end{Verbatim}
\end{itemize}

\textbf{Примери:}
\begin{itemize}
    \item \textbf{Улаз 1:} $1 \rightarrow 4 \rightarrow 5$ \\
          \textbf{Излаз 1:} \texttt{true}
    \item \textbf{Улаз 2:} $1 \rightarrow 2 \rightarrow 5$ \\
          \textbf{Излаз 2:} \texttt{false}
\end{itemize}

\sablon{Шаблон првог теста новембар 2023. група E}{2023/Novembar/GrupaE/AISP2023_Test1_GrupaE.linq}
